\clearpage
\section*{2. Nhận diện biểu cảm khuôn mặt}
\addcontentsline{toc}{section}{2. Nhận diện biểu cảm khuôn mặt}

\subsection*{2.1. Mô hình FER và nhãn cảm xúc}

Hệ thống sử dụng thư viện FER \cite{ferlibrary}, dựa trên bài toán nhận diện
biểu cảm FER2013 \cite{fer2013}. Detector được khởi tạo bằng:

\begin{svcode}
FER(use_tflite=True, mtcnn=False)
\end{svcode}

Chế độ TensorFlow Lite \cite{tensorflowlite} giảm phần phụ thuộc và phù hợp với
inference CPU. MTCNN không cần thiết cho bài này; detector Haar cascade mặc
định của FER được dùng để tìm khuôn mặt. Bảy nhãn thô của mô hình gồm
\texttt{angry}, \texttt{disgust}, \texttt{fear}, \texttt{happy},
\texttt{sad}, \texttt{surprise} và \texttt{neutral}. Trong đó
\texttt{disgust} không có hành động tương ứng nên được chuẩn hóa thành trạng
thái không an toàn và yêu cầu dừng.

Adapter \texttt{FerEmotionModel} nhận output công khai của FER, kiểm tra box
$(x,y,w,h)$, kiểm tra toàn bộ score hữu hạn trong $[0,1]$, rồi chọn nhãn có
score lớn nhất. FER thực tế có thể trả bounding box OpenCV dưới dạng
\texttt{numpy.ndarray}; adapter chuyển từng phần tử nguyên về tuple Python
trước khi tạo \texttt{DetectedFace}. Trường hợp này đã được bổ sung regression
test sau lần chạy webcam thật.

\subsection*{2.2. Pipeline webcam theo thời gian thực}

Node nhận diện mở camera bằng OpenCV \cite{opencv} với backend V4L2, mặc định
là \path{/dev/video0}, độ phân giải yêu cầu $640\times480$ và buffer size bằng
1. Frame được mirror trước khi inference và hiển thị để thao tác giống gương.

\begin{enumerate}[leftmargin=1cm,itemsep=0.2em]
  \item Timer capture chạy ở 20~Hz và đọc frame mới nhất từ webcam.
  \item \texttt{InferenceScheduler} cho phép inference ở 5~Hz; sequence chỉ
        tăng sau khi một lần inference hoàn thành.
  \item FER trả danh sách khuôn mặt và score; adapter tạo danh sách
        \texttt{DetectedFace} đã xác thực.
  \item Nếu đúng một khuôn mặt, observation mang nhãn và confidence của khuôn
        mặt đó; controller xử lý ngay khi confidence lớn hơn 0,60. Nếu không
        có hoặc có nhiều khuôn mặt, observation mang
        \texttt{unknown}, confidence 0 và đúng \texttt{face\_count}.
  \item Observation mới được publish lên \path{/emotion_observation}; frame
        vẫn có thể hiển thị ở 20~Hz giữa hai lần inference.
\end{enumerate}

Cách tách capture và inference tránh tăng sequence bằng các frame hiển thị lặp
lại. Nếu FER bị treo, node không thể phát observation mới; watchdog 0,75~s ở
controller sẽ đưa vận tốc về không thay vì giữ lệnh cũ.

\subsection*{2.3. Overlay và phản hồi trạng thái robot}

Trên bản sao của frame, node vẽ box đã clip theo kích thước ảnh, nhãn FER,
confidence, số khuôn mặt và observation mới nhất. Dòng trạng thái robot lấy từ
\path{/emotion_behavior_status}, hiển thị state, biểu cảm đã ổn định, hành động,
reason và phần trăm tiến độ. Frame đầu vào không bị sửa và không có code path
ghi ảnh hoặc video xuống đĩa.

\reportfigure{02_overlay_nhan_dien.png}{Cửa sổ OpenCV hiển thị một khuôn mặt, nhãn FER, confidence, stable emotion và trạng thái robot.}
\reportfigure{03_sau_bieu_cam.png}{Bộ minh chứng sáu biểu cảm happy, sad, neutral, angry, surprise và fear trong điều kiện ánh sáng thử nghiệm.}
\reportfigure{04_topics_observation_status.png}{Hai topic JSON /emotion\_observation và /emotion\_behavior\_status trong một phiên chạy.}

\subsection*{2.4. Môi trường AI độc lập}

TensorFlow, FER và OpenCV pip không được cài vào Python hệ thống của ROS.
Script \path{scripts/setup_face_venv.sh} tạo \path{.venv-face} với
\texttt{--system-site-packages} để vẫn nhìn thấy \texttt{rclpy}, sau đó cài
các phiên bản được pin: NumPy~1.26.4, OpenCV contrib~4.10.0.84,
Pillow~10.4.0, Requests~2.32.3, setuptools~70.3.0 và
TensorFlow CPU~2.15.1. FER~25.10.3 được cài bằng \texttt{--no-deps} để tránh
kéo nhánh MTCNN/PyTorch không sử dụng.

Launch file chỉ prefix node nhận diện bằng Python của virtualenv. Controller
vẫn chạy bằng ROS Python bình thường. Biên giao tiếp giữa hai môi trường chỉ
gồm message chuẩn, vì vậy dependency AI không can thiệp vào vòng điều khiển.

\subsection*{2.5. Xử lý lỗi và quyền riêng tư}

Khi camera mất frame, node phát một observation \texttt{unknown} mới nếu
publisher còn hoạt động rồi thoát lỗi; controller có cả tín hiệu dừng tức thời
và watchdog dự phòng. Nhấn \texttt{Q}, Ctrl+C hoặc ROS shutdown sẽ release
camera và đóng cửa sổ. Lần thử đầu tiên trên Lenovo Legion 5 cho thấy camera
bị tắt bằng E-shutter nên không tồn tại \path{/dev/video0}; sau khi bật camera,
node mở webcam thành công. Một smoke test headless 12~giây đã xác nhận FER xử
lý frame thật mà không còn crash do kiểu bounding box.

Bài thực hành chỉ suy luận biểu cảm; không nhận diện danh tính, tuổi, giới tính
hoặc chủng tộc. Hệ thống không tự lưu frame, crop mặt hay video. Ảnh minh chứng
có khuôn mặt chỉ được thêm khi người thực hiện chủ động chọn và kiểm tra nội
dung.
